% Part: lambda-calculus
% Chapter: syntax
% Section: unique-readability

\documentclass[../../../include/open-logic-section]{subfiles}

\begin{document}

\olfileid{lam}{syn}{unq}
\olsection{Cara Maca Tunggal}

Kita bisa takon apa saben suku mung nduweni
siji cara kanggo mbentuk, lan jawabane ya.
Saben suku lambda mung nduweni siji cara
pambentukan lan siji cara panapsiran.
Ing rembugan ing ngisor iki,
\emph{pambentukan} iku tata cara mbangun
suku kanthi ngetrapake paugeran pambentukan
(sepisan utawa luwih) saka
\olref[trm]{defn:term}.

\begin{lem}\ollabel{lem:term-start}
Suku diwiwiti dening variabel utawa
tandha kurung bukak.
\end{lem}

\begin{proof}
Sawijining ekspresi diarani suku mung yen
dibangun miturut \olref[trm]{defn:term}.
Yen ekspresi iku asil saka
\olref[trm]{defn:term-var}, mesthi
awujud variabel. Yen iku asil saka
\olref[trm]{defn:term-abs} utawa
\olref[trm]{defn:term-app}, ekspresi
iku diwiwiti dening tandha kurung bukak.
\end{proof}

\begin{lem}\ollabel{lem:app-start}
Asil aplikasi diwiwiti dening rong
tandha kurung bukak utawa dening
tandha kurung bukak banjur variabel.
\end{lem}

\begin{proof}
Yen $M$ iku asil aplikasi, wujude
$(PQ)$, mula diwiwiti dening tandha
kurung bukak. Amarga $P$ iku suku,
miturut \olref{lem:term-start},
suku iku diwiwiti dening tandha kurung
bukak utawa variabel.
\end{proof}

\begin{lem}\ollabel{lem:initial}
Ora ana perangan wiwitan sejati
saka sawijining suku kang dhewe
uga suku.
\end{lem}

\begin{prob}
Buktikna \olref[lam][syn][unq]{lem:initial}
kanthi induksi ing dawa suku.
\end{prob}

\begin{prop}[Cara Maca Tunggal] \ollabel{prop:unq}
Saben suku nduweni pambentukan kang
mung siji. Tegese, yen suku $M$
dibentuk liwat sawijining pambentukan,
mung pambentukan iku kang bisa
mbentuk suku kasebut.
\end{prop}

\begin{proof}
  Kita mbuktekake kanthi induksi ing
  pambentukan suku.
  \begin{enumerate}
    \item $M$ awujud
      $x$, ing ngendi $x$ iku sawijining
      variabel. Amarga asil abstraksi
      lan aplikasi tansah diwiwiti
      dening tandha kurung bukak, rong
      paugeran iku ora bisa digunakake
      kanggo mbangun $M$. Mula
      pambentukan $M$ mesthi mung
      siji langkah saka
      \olref[trm]{defn:term},
      yaiku \olref[trm]{defn:term-var}.
    \item $M$ awujud
      $(\lambd[x][N])$, ing ngendi
      $x$ iku sawijining variabel lan
      $N$ iku suku. Suku iki ora bisa
      dibangun miturut
      \olref[trm]{defn:term}
      \olref[trm]{defn:term-var},
      amarga dudu variabel tunggal.
      Miturut \olref{lem:app-start},
      iku dudu asil aplikasi. Dadi
      $M$ mung bisa dadi asil abstraksi
      saka~$N$. Miturut hipotesis
      induksi, pambentukan $N$ uga
      mung siji.
    \item $M$ awujud $(PQ)$, ing
      ngendi $P$ lan $Q$ iku suku.
      Amarga diwiwiti dening tandha
      kurung bukak, suku iki ora bisa
      dibangun miturut
      \olref[trm]{defn:term}
      \olref[trm]{defn:term-var}.
      Miturut \olref{lem:term-start},
      $P$ ora bisa diwiwiti dening
      $\lambd$, mula $(PQ)$ ora
      bisa dadi asil abstraksi.
      Saiki upamana ana cara liyane
      kanggo mbangun $M$ kanthi
      aplikasi, umpamane uga awujud
      $(P'Q')$. Yen pamisahan iki
      beda, $P$ dadi perangan wiwitan
      sejati saka $P'$ utawa kosokbaline;
      iki ora mungkin miturut
      \olref{lem:initial}. Mula
      $P$ lan $Q$ ditemtokake kanthi
      unik, lan miturut hipotesis
      induksi, pambentukan $P$
      lan $Q$ uga mung siji.
  \end{enumerate}
\end{proof}

Andharan kang luwih gampang diwaca
saka proposisi ing ndhuwur kaya
mangkene:
\begin{prop}
  Suku $M$ mung bisa awujud salah
  siji saka wujud iki:
  \begin{enumerate}
    \item $x$, ing ngendi $x$ iku
      variabel kang ditemtokake kanthi
      unik dening $M$.
    \item $(\lambd[x][N])$, ing ngendi
      $x$ iku variabel lan $N$ iku
      suku liyane; loro-lorone
      ditemtokake kanthi unik
      dening $M$.
    \item $(PQ)$, ing ngendi $P$ lan
      $Q$ iku rong suku kang
      ditemtokake kanthi unik
      dening $M$.
  \end{enumerate}
\end{prop}

\end{document}
